% ============================================
% Johns Hopkins Letter Template
% Assistant Professor Cover Letter – Arizona State University (W. P. Carey)
% ============================================

\documentclass[11pt,letterpaper]{letter}

\usepackage{graphicx}
\usepackage{eso-pic}
\usepackage{geometry}
\usepackage{microtype}
\usepackage[hidelinks]{hyperref}

% ----- Page Setup -----
\geometry{left=25mm,right=25mm,top=25mm,bottom=25mm}

% ----- Logo Placement (first page only) -----
\newcommand{\PutJHULogoFG}[1]{%
  \AddToShipoutPictureFG*{%
    \AtPageUpperLeft{%
      \hspace{10mm}%
      \raisebox{-38mm}{%
        \includegraphics[width=6cm]{#1}%
      }%
    }%
  }%
}

% ----- Sender Information -----
\signature{Decory Edwards}
\address{Department of Economics \\ Johns Hopkins University \\ Baltimore, MD 21218 \\ \texttt{dedwar65@jh.edu}}

\begin{document}
\PutJHULogoFG{Images/jhulogo.png}

\begin{letter}{Search Committee \\
Department of Economics \\
W. P. Carey School of Business \\
Arizona State University \\
Tempe, AZ 85287 \\ USA}

\opening{Dear Members of the Search Committee,}

I am writing to express my interest in the Assistant Professor position in the Department of Economics at the W. P. Carey School of Business at Arizona State University, beginning Fall 2026. I am a Ph.D. candidate in Economics at Johns Hopkins University, advised by Professors Christopher D.~Carroll, Nicholas W.~Papageorge, and Stelios Fourakis, and I expect to receive my degree in May 2026. My research fields are macroeconomics, wealth and income dynamics, and computational economics.

My research agenda focuses on understanding the determinants of wealth inequality and the mechanisms through which heterogeneity in returns to wealth shapes aggregate outcomes. In my job-market paper, I estimate the degree of heterogeneity in individual portfolio returns using micro-data from the Survey of Consumer Finances and simulate a structural model in which households face idiosyncratic return risk. I show that allowing for persistent heterogeneity in returns substantially improves the model’s ability to match the empirical distribution of wealth, offering a tractable way to connect micro-level return dispersion to macroeconomic inequality.

In a second project, I use the Health and Retirement Study to examine the relationship between interpersonal trust and portfolio performance. Building on the literature that finds a hump-shaped relationship between trust and income, I show that a similar relationship holds for individual returns to wealth measured in the HRS data, highlighting a behavioral channel through which social attitudes shape saving and investment outcomes. This work also provides new estimates of the persistent component of returns to net worth, which motivate the calibration of heterogeneity in my structural model.

ASU’s Department of Economics is an internationally recognized research department whose faculty regularly publish in the top journals in macroeconomics, applied microeconomics, and economic theory. The position’s emphasis on an ambitious research agenda, high-quality instruction at both the undergraduate and graduate levels, and active engagement in mentoring doctoral students strongly aligns with my training and research goals. I am also drawn to the School’s commitment to inclusive excellence and to ASU’s Charter—particularly its focus on broad access, public value, and serving diverse communities. My work in computational macroeconomics and heterogeneous-agent modeling fits well within a research environment that values innovation, quantitative rigor, and the effective integration of AI-related tools into research and teaching.

\textbf{Teaching.} My teaching experience centers on undergraduate macroeconomics and an upper-level macro policy seminar, where I have led weekly discussion sections and held regular office hours for classes ranging from 16 to 40 students. My approach is student-centered: I aim to help students “convince themselves” that they have moved from confusion to understanding by holding structured office-hour coaching that builds trust and confidence. At ASU, I am prepared to teach \textit{Principles of Macroeconomics}, \textit{Intermediate Macroeconomic Theory}, \textit{Money and Banking}, and \textit{Econometrics}, and I look forward to contributing to graduate instruction and doctoral mentoring as well.

I believe I would be a strong fit because my computational and empirical toolkit allows me to (i) produce rigorous, data-backed estimates of returns and distributional outcomes, (ii) build and validate quantitative models that speak to macroeconomic, financial, and policy-relevant questions, and (iii) mentor students effectively at both the undergraduate and graduate levels. I am enthusiastic about the opportunity to contribute to ASU’s teaching, research, and service missions within the W. P. Carey School of Business.

I would be honored to join Arizona State University’s Department of Economics. Thank you for your consideration; I welcome the opportunity to discuss my fit with your department at your convenience.

\closing{Sincerely,}

\end{letter}
\end{document}
